\section{An intrinsic hierarchy for paid causal synthesis}\label{sec:synthesis}
A comparison inequality does not create the acquisition program used by its upper bound. We give an intrinsic criterion over all classical feedback programs on a finite-branching causal interface. It applies to unknown-model classical experiments and finite-dimensional quantum instruments, with one program shared by every model and initial state. Exact unrestricted-history synthesis and exact fixed-memory synthesis remain different conclusions.

\subsection{Primitive data and the controller class}
A visible history records executed commands and reports. At each open decision node $h$, the legal next-command set $\mathcal A(h)$ is finite and nonempty. After a productive report there may also be a finite set $\mathcal O(h)$ of legal terminal routes. There is no terminal route at the root or after a nonproductive report. Each command has finitely many reports, so the tree is locally finite and countable. Legal domains are common to the admitted models and are part of the interface. A physical phase determined by visible history is not an additional free input to a finite-memory observer.

The unknown model $\theta$ belongs to an arbitrary nonempty set $\Theta$, not necessarily finite, compact or equipped with a prior. At each history the hidden quantum state lies in a specified finite-dimensional matrix algebra. Command $a$ has completely positive maps $\mathcal I_{\theta,h,a,y}$ with
\[
 \sum_y\mathcal I_{\theta,h,a,y}^*(I)=I.
\]
Dimensions may depend on the model and finite history. The initial identity is $I_\theta$. Nonnegative classical transition matrices are the diagonal special case. There is no positive-likelihood or eigenvalue floor.

Costs are nonnegative. Their conditional expected-cost observables satisfy the uniform floor
\begin{equation}\label{eq:syn-cost-floor}
 G_{\theta,h,a}\succeq\underline c I,\qquad \underline c>0.
\end{equation}
A pathwise floor implies this assumption, but only the displayed conditional floor is needed here. The cost $C_\tau$ is the core cost through commitment. A separately executed preparation or audit is charged in addition or represented as a compulsory call; it is not silently included without its physical operation.

Fix finitely many bounded terminal scores $j=1,\ldots,J$. Route $o$ at $h$ has a Hermitian score observable $F_{\theta,h,o,j}$ on the current state, with
\begin{equation}\label{eq:syn-score-bounds}
 \alpha_jI\preccurlyeq F_{\theta,h,o,j}\preccurlyeq\beta_jI,
 \qquad -\infty<\alpha_j\le\beta_j<\infty.
\end{equation}
The route is chosen before observing its future audit. For example, $F$ is the adjoint of the declared audit instrument applied to a bounded score of its outcome and the chosen route. These are primitive unnormalized observables, not assumed posteriors. Prescribed Hermitian targets on the initial state are denoted by $A_{\theta,j}$; constant targets are scalar multiples of $I_\theta$.

A controller uses full visible history and independent private randomness, but neither the unknown model nor an unknown initial density matrix. It must commit exactly once almost surely and satisfy
\begin{equation}\label{eq:syn-goal}
 \sup_{\theta,\rho}\E_{\theta,\rho}C_\tau\le N,\qquad
 \E_{\theta,\rho}\Phi_j=\tr(\rho A_{\theta,j})
 \quad(\theta\in\Theta,\ \rho\succeq0,\ \tr\rho=1).
\end{equation}
This is classical feedback over the specified quantum experiment, not an uncharged coherent controller memory. Uniformity over initial states is distinct from uniformity over rooted continuations of a history-dependent experiment; both are required before a program may be repeated after rejection. Restricted preparation faces and deadlocks are treated explicitly in the current supplement.

\subsection{One realization plan, rather than modelwise policies}
At a horizon $T\ge1$, include decisions following the first $T$ reports. Continuing at depth $T$ leaves an open branch; its next physical call has not occurred. The weight $\gamma(h)$ is the product of controller choice probabilities along $h$, excluding physical report probabilities. Write $\gamma(h,d)$ after a command or terminal decision. The common sequence-form polytope $\mathcal P_T$ is given by
\begin{equation}\label{eq:realization-plan}
 \gamma(\varnothing)=1,\quad \gamma(h,d)\ge0,\quad
 \sum_d\gamma(h,d)=\gamma(h),\quad
 \gamma(h,a,y)=\gamma(h,a)
\end{equation}
for every present report child. All variables lie in $[0,1]$, and the same variables are used for every model. Sequence form itself is classical \cite{Stengel96}; its role is to preserve the common-policy quantifier.

Let $\mathcal I_{\theta,h}$ be the composition of the primitive maps along $h$, without controller weights. Define affine functions on this polytope by
\begin{align}
 \mathsf B_{\theta,T}&=\sum_{\substack{|h|\le T\\o\in\mathcal O(h)}}
       \gamma(h,o)\mathcal I_{\theta,h}^*(I),&
 \mathsf S_{\theta,T}&=I_\theta-\mathsf B_{\theta,T},\label{eq:syn-mass}\\
 \mathsf V_{\theta,j,T}&=\sum_{\substack{|h|\le T\\o\in\mathcal O(h)}}
       \gamma(h,o)\mathcal I_{\theta,h}^*(F_{\theta,h,o,j}),\label{eq:syn-value}\\
 \mathsf C_{\theta,T}&=\sum_{\substack{|h|<T\\a\in\mathcal A(h)}}
       \gamma(h,a)\mathcal I_{\theta,h}^*(G_{\theta,h,a}).\label{eq:syn-cost}
\end{align}
Every coefficient is a raw finite-path matrix. No report probability is divided out.

\begin{lemma}[One plan for all models]\label{lem:common-plan}
The polytope is nonempty and compact. Its elements are exactly the finite-prefix laws of common randomized feedback controllers, with expectations given by traces against the displayed operators, simultaneously for every model and initial state. Arbitrary private random memory does not enlarge these laws relative to visible-history policies using fresh randomness. Moreover $\mathsf B_{\theta,T},\mathsf S_{\theta,T}\succeq0$ and their sum is $I_\theta$.
\end{lemma}
\begin{proof}
A behavioral policy gives the product weights. Conversely, at a node of positive weight choose $d$ with probability $\gamma(h,d)/\gamma(h)$; at a zero-weight node choose any legal row. Induction recovers every specified weight. The finite linear constraints form a closed bounded set, and any choice of legal rows proves nonemptiness.

To include arbitrary private memory, fix a visible action--report path and integrate over all private random tapes producing those controller choices when supplied that report path. This probability is independent of the physical model and initial state. The raw path instrument is the same for every compatible tape, because the executed commands and reports have been fixed. Thus it factors from the tape integral, and the integrated choice weights obey the same equations. The ratio construction reproduces the laws in all models at once. It uses full visible-history storage; it does not make a tape into free finite memory.

Multiply the raw path effects by these weights. Summing report effects uses normalization of each instrument, and summing controller choices uses the plan equation. Induction shows that terminated and open effects sum to the incoming effect, starting from $I_\theta$. Every such effect is positive. The cost and score formulas follow by summing their actual finite-prefix contributions.
\end{proof}

\subsection{The necessary cost reserve and terminal-score intervals}
For a finite $F\subset\Theta$, let $\mathcal H_T(F,N)$ be the feasibility problem in $\gamma\in\mathcal P_T$ defined by
\begin{align}
 \mathsf C_{\theta,T}+\underline c\mathsf S_{\theta,T}
       &\preccurlyeq NI_\theta,\label{eq:hier-cost}\\
 \alpha_j\mathsf S_{\theta,T}
       &\preccurlyeq A_{\theta,j}-\mathsf V_{\theta,j,T}
        \preccurlyeq\beta_j\mathsf S_{\theta,T}
       \quad(\theta\in F,\ 1\le j\le J).\label{eq:hier-tail}
\end{align}
These are linear inequalities in the classical case and linear matrix inequalities in the quantum case. They are not implemented by pretending that open paths terminate at the cutoff. Those paths owe at least one further paid call, and their remaining bounded scores lie in the displayed interval.

\begin{lemma}[Restriction and the sharp budget tail]\label{lem:synthesis-tail}
A feasible plan at level $T+1$ restricts to a feasible plan at level $T$. Every feasible level with $N\ge\underline c$ satisfies
\begin{equation}\label{eq:sharp-tail}
 0\preccurlyeq\mathsf S_{\theta,T}\preccurlyeq
 \delta_T I_\theta,\qquad
 \delta_T=\min\{1,(N/\underline c-1)/T\}.
\end{equation}
If $N<\underline c$, there is no feasible proper program.
\end{lemma}
\begin{proof}
Write $D_T=\mathsf B_{T+1}-\mathsf B_T\succeq0$ for the new terminal effect. Normalization gives $\mathsf S_T=D_T+\mathsf S_{T+1}$. New score increments lie between $\alpha_jD_T$ and $\beta_jD_T$. Adding them to the next-level remainder intervals proves the current-level intervals. Every path still open after depth $T$ executes another call before the next decision. Therefore $\mathsf C_{T+1}-\mathsf C_T\succeq\underline c\mathsf S_T$. Together with the next-level cost constraint this implies the current one.

Let $S_t$ denote the open effect after $t$ calls, so $S_0=I$ and $S_t\succeq S_T$ for $t<T$. The conditional cost floor gives $\mathsf C_T\succeq\underline c\sum_{t=0}^{T-1}S_t\succeq\underline c[I+(T-1)S_T]$. Adding the reserve gives $\underline c[I+TS_T]\preccurlyeq NI$. Positivity and $S_T\preccurlyeq I$ prove the claimed bound. Every proper program must execute its first paid call, proving the infeasible boundary below the floor.
\end{proof}

\begin{theorem}[Finite-obstruction characterization of paid causal synthesis]\label{thm:causal-hierarchy}
Under the stated finite-branching, positive-cost and bounded-score hypotheses, the following are equivalent:

\textup{(i)} One common proper visible-history program satisfies \eqref{eq:syn-goal} for every model and every initial state.

\textup{(ii)} $\mathcal H_T(F,N)$ is feasible for every finite model set $F\subset\Theta$ and every integer $T\ge1$.

Consequently nonexistence is witnessed by a finite horizon and a finite set of operating models. Feasibility at arbitrarily many tested levels alone is not a finite certificate for the infinite condition. No fixed bound on controller memory is asserted for the limiting policy.
\end{theorem}
\begin{proof}
For necessity take the prefix of the actual program. Every open path must make at least one additional call before commitment, so its conditional future expected cost is at least $\underline c$. Nonnegative conditional summation proves the reserved cost inequality. Scores on the disjoint future terminal events lie between $\alpha_j$ and $\beta_j$ times their total open mass. Almost-sure commitment makes that mass exactly $\mathsf S_T$, proving the score intervals. These same prefix weights work in every model.

For sufficiency consider the compact product of all nodewise probability simplexes on the countable history tree. A finite-prefix plan can be extended to an element of this product by arbitrary legal rows at later nodes. Its realization weights are finite products of coordinates, hence continuous. Each finite collection of the cost and score constraints is closed. For such a collection take the union of its finitely many model indices and its greatest horizon. Assumption (ii) and the restriction lemma give one feasible plan for every constraint in the collection. Extend it to all nodes. The closed constraints therefore have the finite-intersection property in a compact space, and one common infinite policy satisfies them all.

For this policy \eqref{eq:sharp-tail} gives uniformly vanishing survival probabilities, including the exact boundary $N=\underline c$. Thus it commits almost surely in every model and initial state. Its partial costs are nonnegative and at most $N$; monotone convergence yields the actual total cost bound. The norm of the score remainder is at most $\max(|\alpha_j|,|\beta_j|)\delta_T$. It tends to zero, so the bounded terminal scores have exactly the prescribed operator expectations. The program is common because the product coordinates were common throughout. Finally, negating the equivalence proves existence of a finite obstruction when no such program exists.
\end{proof}

\begin{corollary}[Attainment of the least finite cost]\label{cor:synthesis-attainment}
If any common proper program satisfies the prescribed targets at finite uniform cost, then the infimum of such cost budgets is attained.
\end{corollary}
\begin{proof}
Let feasible budgets decrease to their finite infimum $N_*$. For each fixed finite model set and horizon, the corresponding nonempty compact feasible sets decrease with the budget. Their intersection is nonempty and is feasible at $N_*$. Every finite problem at $N_*$ is therefore feasible; the theorem gives a proper program at that budget. The first-call floor ensures $N_*\ge\underline c$.
\end{proof}

\subsection{A strict dual certificate at an obstructing level}
List all left-minus-right matrix violations of \eqref{eq:hier-cost}--\eqref{eq:hier-tail} as $K_l(\gamma)$, $1\le l\le L_0$. Thus feasibility is $K_l\preccurlyeq0$ for every $l$. Define
\[
 v=\min_{\gamma\in\mathcal P_T}
       \max\{0,\max_l\lambda_{\max}(K_l(\gamma))\}.
\]
Compactness makes this minimum attained and $v>0$ exactly when the finite problem is infeasible.

\begin{theorem}[Strict finite dual obstruction]\label{thm:synthesis-dual}
The exact violation value is
\begin{equation}\label{eq:synthesis-dual}
 v=\max_{\substack{W_l\succeq0\\\sum_l\tr W_l\le1}}
       \min_{\gamma\in\mathcal P_T}\sum_l\tr[W_lK_l(\gamma)].
\end{equation}
The maximum is attained. An infeasible level has a positive-matrix separating family with strictly positive value. The inner minimization is a finite linear program in the common realization weights. No Slater hypothesis or inverse covariance is required.
\end{theorem}
\begin{proof}
For a fixed collection of violations, maximizing the displayed pairing over positive matrices with total trace at most one gives their largest positive eigenvalue, or zero by taking all matrices zero. Taking a minimizer defining $v$ gives the weak upper bound for every family $W$.

For the converse when $v>0$, fix $0\le\eta<v$. In the finite direct sum of Hermitian matrix spaces, the point $(\eta I_l)_l$ does not belong to the closed convex set
\[
 \{(K_l(\gamma))_l:\gamma\in\mathcal P_T\}
       +\prod_l\{D_l\succeq0\}.
\]
This set is closed because its first summand is compact. Finite-dimensional strict separation gives a nonzero linear functional whose infimum on that set exceeds its value at the excluded point. The functional must be nonnegative on the positive cone, or its infimum along a cone ray would be minus infinity. Hence it has the form of pairing with positive $W_l$. Their total trace is positive; normalize it to one. Separation yields an inner minimum greater than $\eta$. Let $\eta\uparrow v$ and take a convergent subsequence in the compact normalized positive-matrix set. The inner minimum is continuous in $W$, being the minimum of a continuous bilinear function over compact $\mathcal P_T$. Its limiting value is at least $v$. If $v=0$, a feasible plan gives the upper bound zero and $W=0$ attains zero. This proves equality and attainment in all cases.
\end{proof}

The finite dual is a genuine obstruction, not an assertion that the first obstructing horizon is uniformly bounded or effectively discoverable for arbitrary real primitive data. Existing quantum decision-process and reachability theories already contain substantial control, duality and computability results \cite{Barry14,Ying18,Saldi25}. Generic quantum MDPs, sequence form and convex separation are not claimed as new ingredients.

\subsection{Paid finite-prefix completion}
Suppose every open leaf reachable by a selected finite prefix admits one common legal fallback that commits with conditional expected cost at most $\kappa_0\ge\underline c$, uses at most $B_0$ live roles, and has terminal scores in the same $[\alpha_j,\beta_j]$ intervals. Its command and physical phase compatibility must be certified at these leaves.

\begin{corollary}[A paid finite-prefix completion]\label{cor:synthesis-fallback}
A feasible level-$T$ plan completed by this fallback has, for every model and initial state,
\begin{align}
 \E C&\le N+(\kappa_0-\underline c)\delta_T,\label{eq:fallback-budget}\\
 - (\beta_j-\alpha_j)\delta_T I
 &\preccurlyeq A_{\theta,j}-\E[\Phi_j]\preccurlyeq
       (\beta_j-\alpha_j)\delta_T I.\label{eq:fallback-error}
\end{align}
Here the score expectation is its operator on the initial state. If the finite prefix uses $L_T$ live node roles and the completed terminal decoder has $t$ actual values, a sufficient state-stopped label count is $L_T+B_0+t$ and an event-stopped count is $\max(L_T+B_0,t)$. Program constants, arithmetic and the possible exponential growth of $L_T$ are separate costs.
\end{corollary}
\begin{proof}
The completed cost is at most $\mathsf C_T+\kappa_0\mathsf S_T$. Subtract the reserved cost inequality and use the survival bound. Both the prescribed remaining score and the fallback's remaining score lie between $\alpha_j\mathsf S_T$ and $\beta_j\mathsf S_T$. Their difference lies between the stated opposite multiples of $\mathsf S_T$, proving the error bound. A finite prefix is implemented by storing its current node; the fallback uses its declared live roles. The exact role-packing lemma gives the two sufficient counts. These are existence counts, not a claim that a real-valued finite plan has an exact finite-bit sampler.
\end{proof}

\begin{proposition}[Sharpness of the uniform tail rate]\label{prop:tail-sharp}
For unit cost per call, every $b\ge1$ and integer $T\ge1$ with $b-1\le T$ admit a proper finite-branching program of mean cost $b$ and survival probability $(b-1)/T$ after $T$ calls.
\end{proposition}
\begin{proof}
With probability $p=(b-1)/T$ commit on call $T+1$, and otherwise on call one. Then $\E\tau=1+Tp=b$ and $\Pr(\tau>T)=p$. Its prefix has cost $1+(T-1)p$ and reserve $p$. Thus the reserve and the boundary $b=1$ are exact. If $b-1>T$, the universal tail bound one is attained by a deterministic sufficiently long path. This concerns the survival bound, not a uniform expected cost-tail estimate.
\end{proof}

\subsection{Synthesis of a finite marked squared-task reference}
For a finite virtual mark set $X$, fix masses $\nu_x$, target means $f(x)\in\R^k$ and nonnegative residuals $w(x)$ of a genuinely prepared reference law. Assume bounded audit coordinates. Use as terminal scores
\[
 \mathbf1_{\{X=x\}},\quad
 \mathbf1_{\{X=x\}}Z_i\ (1\le i\le k),\quad
 \mathbf1_{\{X=x\}}\|Z\|^2,
\]
with respective scalar-identity targets $\nu_x$, $\nu_x f_i(x)$ and $\nu_x[w(x)+\|f(x)\|^2]$. They are computed from the declared terminal audit instrument and route.

\begin{corollary}[Intrinsic synthesis of finite-mark squared-task balance]\label{cor:balance-hierarchy}
There is one common proper finite-branching acquisition program with these unnormalized marked moments and cost at most $N$ for all models and initial states if and only if every finite problem of the hierarchy for this finite score list is feasible. Such a program satisfies
\[
 \E[q(X)\|Z-r(X)\|^2]
   =\sum_x\nu_x q(x)[w(x)+\|f(x)-r(x)\|^2]
\]
for every selected readout. If, in addition, the same return rule is legal and the same certificate holds on every actual rooted continuation after rejection, it gives exact task balance for repeated acquisition. The hierarchy alone does not supply a finite live-memory bound.
\end{corollary}
\begin{proof}
The first equivalence is the hierarchy theorem applied to the explicit bounded score list. Expand the squared loss in each mark and substitute the three unnormalized moment identities. This proves the task equation without division by $\nu_x$, so zero-mass marks are harmless. Conversely the displayed marked moment requirement is exactly that finite target list. Applying this program repeatedly requires its hypotheses on the actual later roots; no equality of initial-state matrices can replace missing history-dependent continuation data. A finite prefix with a legal fallback, or the separate finite-memory theorem, supplies an explicitly charged implementation when its additional hypotheses hold.
\end{proof}
